% =============================================================================
% Consolidated Research Foundations for AG-SAR
% Hallucination Detection in Large Language Models: Theory and Methods
% =============================================================================
% This document consolidates the foundational research from six key papers:
% 1. V-Information (Xu et al., ICLR 2020) - Bounded information theory
% 2. Attention Rollout (Abnar & Zuidema, ACL 2020) - Attention flow analysis
% 3. Semantic Entropy (Kuhn et al., ICLR 2023) - Uncertainty estimation
% 4. SelfCheckGPT (Manakul et al., EMNLP 2023) - Black-box hallucination detection
% 5. DoLa (Chuang et al., ICLR 2024) - Layer contrasting for factuality
% 6. Epiplexity (Finzi et al., JMLR 2025) - Computational information bounds
% =============================================================================

\documentclass[11pt]{article}

% Essential packages
\usepackage[utf8]{inputenc}
\usepackage[T1]{fontenc}
\usepackage{amsmath,amssymb,amsthm}
\usepackage{mathtools}
\usepackage{hyperref}

% Theorem environments
\newtheorem{theorem}{Theorem}[section]
\newtheorem{lemma}[theorem]{Lemma}
\newtheorem{proposition}[theorem]{Proposition}
\newtheorem{corollary}[theorem]{Corollary}
\newtheorem{definition}[theorem]{Definition}
\newtheorem{assumption}[theorem]{Assumption}
\newtheorem{remark}[theorem]{Remark}

% Common notation
\newcommand{\E}{\mathbb{E}}
\newcommand{\R}{\mathbb{R}}
\newcommand{\N}{\mathcal{N}}
\newcommand{\X}{\mathcal{X}}
\newcommand{\Y}{\mathcal{Y}}
\newcommand{\F}{\mathcal{F}}
\newcommand{\D}{\mathcal{D}}
\newcommand{\V}{\mathcal{V}}
\newcommand{\KL}{\mathrm{KL}}
\newcommand{\JSD}{\mathrm{JSD}}
\newcommand{\Var}{\mathrm{Var}}
\newcommand{\Cov}{\mathrm{Cov}}
\newcommand{\tr}{\mathrm{tr}}
\DeclareMathOperator*{\argmax}{arg\,max}
\DeclareMathOperator*{\argmin}{arg\,min}
\DeclareMathOperator{\softmax}{softmax}

\title{Hallucination Detection in Large Language Models:\\
Theoretical Foundations and Methods}

\author{Consolidated Research Summary}

\date{}

\begin{document}

\maketitle

\begin{abstract}
This document consolidates the theoretical foundations and methodological advances underlying modern hallucination detection in large language models (LLMs). We synthesize six foundational works spanning information theory under computational constraints (V-Information, Epiplexity), attention flow analysis (Attention Rollout), uncertainty quantification (Semantic Entropy), and hallucination detection methods (SelfCheckGPT, DoLa). The core insight unifying these works is that hallucination detection requires understanding both the information-theoretic properties of model outputs and the internal mechanisms by which LLMs process and transform information across layers. We present unified notation and highlight the mathematical connections between these approaches, providing a comprehensive reference for researchers developing next-generation hallucination detection systems.
\end{abstract}

\tableofcontents
\newpage

% =============================================================================
\section{Introduction and Motivation}
% =============================================================================

Large language models have demonstrated remarkable capabilities across diverse tasks, yet their tendency to generate content that deviates from factual information---commonly termed ``hallucination''---remains a critical challenge. This problem is particularly acute in high-stakes applications where reliable generation of trustworthy text is crucial.

The hallucination problem can be understood from multiple perspectives:
\begin{enumerate}
    \item \textbf{Information-theoretic}: What information is actually accessible to a computationally bounded model, and how does this differ from the Shannon-theoretic ideal?
    \item \textbf{Mechanistic}: How do transformers process and propagate information through attention mechanisms, and where does factual knowledge localize?
    \item \textbf{Uncertainty quantification}: How can we distinguish model confidence from actual reliability?
    \item \textbf{Detection}: What practical methods can identify hallucinated content at inference time?
\end{enumerate}

This document addresses each perspective, providing a unified theoretical foundation.

% =============================================================================
\section{Information Theory Under Computational Constraints}
% =============================================================================

\subsection{The Limitations of Shannon Information}

Classical information theory assumes observers with unlimited computational capacity. However, this assumption leads to counterintuitive results in machine learning contexts. Consider encrypted messages: according to Shannon theory, encrypted text has high mutual information with the plaintext, yet to any practical observer (or ML model), the encrypted text provides no usable information without the decryption key.

This observation motivates a fundamental question: \emph{What is the usable information content accessible to computationally bounded observers?}

\subsection{Predictive V-Information}

The theory of \emph{predictive $\F$-information} addresses this question by parameterizing information measures with a family of predictive functions $\F$ representing the observer's computational capabilities.

\begin{definition}[Predictive Family]
A predictive family $\F$ is a set of functions mapping inputs to probability distributions over outputs. Each $f \in \F$ maps $\X \cup \{\varnothing\} \to \mathcal{P}(\Y)$, where $f[\varnothing]$ represents the prediction without side information (the prior), and $f[x]$ represents the prediction conditioned on observing $x$.
\end{definition}

\begin{definition}[$\F$-Entropy]
The $\F$-entropy of a random variable $Y$ is:
\begin{equation}
H_\F(Y) = \inf_{f \in \F} \E_{y \sim Y}\left[-\log f[\varnothing](y)\right]
\end{equation}
This represents the minimum expected surprise achievable by any predictor in $\F$ without access to side information.
\end{definition}

\begin{definition}[$\F$-Information]
The predictive $\F$-information from $X$ to $Y$ is:
\begin{equation}
I_\F(X \to Y) = H_\F(Y) - H_\F(Y \mid X)
\end{equation}
where the conditional $\F$-entropy is:
\begin{equation}
H_\F(Y \mid X) = \inf_{f \in \F} \E_{x,y \sim X,Y}\left[-\log f[x](y)\right]
\end{equation}
\end{definition}

\subsubsection{Special Cases}

The framework recovers well-known quantities under specific choices of $\F$:

\begin{proposition}[Unified Framework]
Let $\Omega$ be the set of all measurable functions. Then:
\begin{enumerate}
    \item $H_\Omega(Y)$ equals Shannon entropy
    \item $I_\Omega(X \to Y)$ equals Shannon mutual information
    \item For $\F$ restricted to linear Gaussian predictors with $\Y = \R^d$, $I_\F(X \to Y)$ equals $R^2 \cdot \tr(\Cov(Y))$
    \item For $\F$ restricted to exponential family predictors with sufficient statistics $\mathbf{t}$, $H_\F(Y)$ equals the maximum entropy distribution with matching expected sufficient statistics
\end{enumerate}
\end{proposition}

\subsubsection{Key Properties}

Unlike Shannon information, $\F$-information exhibits properties more aligned with practical intuition:

\begin{proposition}[Properties of $\F$-Information]
\begin{enumerate}
    \item \textbf{Monotonicity}: If $\F \subseteq \mathcal{U}$, then $H_\F(Y) \geq H_\mathcal{U}(Y)$
    \item \textbf{Non-negativity}: $I_\F(X \to Y) \geq 0$
    \item \textbf{Asymmetry}: In general, $I_\F(X \to Y) \neq I_\F(Y \to X)$
    \item \textbf{Data processing can increase information}: For deterministic $t$, it is possible that $I_\F(t(X) \to Y) > I_\F(X \to Y)$
\end{enumerate}
\end{proposition}

The violation of the data processing inequality is crucial: it explains why representation learning and feature extraction can increase ``usable information'' even though Shannon mutual information cannot increase under deterministic transformations.

\subsubsection{PAC Bounds for Estimation}

A key advantage of $\F$-information is tractable estimation with guarantees:

\begin{theorem}[PAC Bound for $\F$-Information]
Let $\D = \{(x_i, y_i)\}_{i=1}^N$ be samples from the joint distribution. Assume $|\log f[x](y)| \leq B$ for all $f \in \F$, $x \in \X$, $y \in \Y$. Then with probability at least $1-2\delta$:
\begin{equation}
\left|I_\F(X \to Y) - \hat{I}_\F(X \to Y; \D)\right| \leq 4\mathfrak{R}_N(\mathcal{G}_\F) + 2B\sqrt{\frac{2\log(1/\delta)}{N}}
\end{equation}
where $\mathfrak{R}_N(\mathcal{G}_\F)$ is the Rademacher complexity of the log-probability function class induced by $\F$.
\end{theorem}

\subsection{Epiplexity: Structural vs Random Information}

While V-information captures usable information, the theory of \emph{epiplexity} further distinguishes between structural (learnable) and random (unpredictable) information components.

\subsubsection{Three Apparent Paradoxes}

Classical information theory produces statements that conflict with empirical observations:

\begin{enumerate}
    \item \textbf{Information cannot be increased by deterministic processes}: Shannon entropy and Kolmogorov complexity cannot increase under deterministic transformations, yet synthetic data improves model capabilities and AlphaZero learns sophisticated strategies through self-play.

    \item \textbf{Information is independent of factorization order}: Shannon entropy is symmetric: $H(X,Y) = H(X) + H(Y|X) = H(Y) + H(X|Y)$. Yet LLMs learn better on left-to-right text than reversed text, and cryptography relies on asymmetric computational difficulty.

    \item \textbf{Likelihood modeling is merely distribution matching}: A model cannot achieve higher likelihood than the true data-generating process, yet models can extract more structure than present in simple generating rules (e.g., emergent behaviors in cellular automata).
\end{enumerate}

\subsubsection{Time-Bounded Entropy and Epiplexity}

The resolution involves distinguishing random from structural information under computational bounds.

\begin{definition}[Time-Bounded Entropy]
For a computational bound $T$ (e.g., polynomial time), the time-bounded entropy $H_T(X)$ is the entropy of $X$ with respect to observers limited to computation time $T$. Specifically, if the observer cannot distinguish $X$ from uniform randomness in time $T$, then $X$ has maximal time-bounded entropy regardless of its Kolmogorov complexity.
\end{definition}

\begin{definition}[Epiplexity]
Epiplexity (epistemic complexity) quantifies the structural information extractable by a computationally bounded observer. Using the minimum description length principle:
\begin{equation}
\text{Epiplexity}(x) = \min_{H \in \mathcal{H}_T} \left[L(H) + H_T(x \mid H)\right] - H_T(x)
\end{equation}
where $\mathcal{H}_T$ is the class of models computable in time $T$, $L(H)$ is the description length of model $H$, and $H_T(x \mid H)$ is the residual time-bounded entropy given the model.
\end{definition}

\subsubsection{Practical Measurement}

Epiplexity can be estimated from neural network training dynamics:
\begin{equation}
\text{Epiplexity} \approx \int_0^\infty \left[\mathcal{L}(t) - \mathcal{L}(\infty)\right] dt
\end{equation}
where $\mathcal{L}(t)$ is the loss at training step $t$. This measures the ``area under the loss curve above the final loss''---the accumulated learning that transfers to model weights.

% =============================================================================
\section{Attention Flow in Transformers}
% =============================================================================

Understanding how information propagates through transformer layers is essential for interpreting model behavior and detecting hallucinations.

\subsection{Attention Mechanisms}

In a transformer layer, attention weights determine how information from different positions is aggregated:
\begin{equation}
\text{Attention}(Q, K, V) = \softmax\left(\frac{QK^\top}{\sqrt{d_k}}\right)V
\end{equation}

For a sequence of length $n$, each layer $\ell$ produces attention weights $A^{(\ell)} \in \R^{n \times n}$, where $A^{(\ell)}_{ij}$ represents how much position $i$ attends to position $j$.

\subsection{Attention Rollout}

Raw attention weights from individual layers provide incomplete information about information flow across the full network. \emph{Attention rollout} addresses this by computing cumulative attention flow.

\begin{definition}[Attention Rollout]
Given attention matrices $A^{(1)}, \ldots, A^{(L)}$ for an $L$-layer transformer, the attention rollout is computed recursively:
\begin{align}
\tilde{A}^{(0)} &= I_n \\
\tilde{A}^{(\ell)} &= A^{(\ell)} \cdot \tilde{A}^{(\ell-1)} \quad \text{for } \ell = 1, \ldots, L
\end{align}
The final rollout matrix $\tilde{A}^{(L)}_{ij}$ quantifies how much the representation at position $i$ in the final layer depends on the input at position $j$.
\end{definition}

\subsection{Accounting for Residual Connections}

Transformers use residual connections, meaning the output of each layer is:
\begin{equation}
h^{(\ell)} = h^{(\ell-1)} + \text{Attention}^{(\ell)}(h^{(\ell-1)})
\end{equation}

To account for this, we modify the attention matrices:
\begin{equation}
\bar{A}^{(\ell)} = 0.5 \cdot I + 0.5 \cdot A^{(\ell)}
\end{equation}

The rollout then becomes:
\begin{equation}
\tilde{A}^{(L)} = \prod_{\ell=1}^{L} \bar{A}^{(\ell)}
\end{equation}

\subsection{Multi-Head Attention}

For multi-head attention with $H$ heads, attention weights must be aggregated. Common approaches:
\begin{enumerate}
    \item \textbf{Mean}: $A^{(\ell)} = \frac{1}{H}\sum_{h=1}^H A^{(\ell)}_h$
    \item \textbf{Max}: $A^{(\ell)}_{ij} = \max_h A^{(\ell)}_{h,ij}$
    \item \textbf{Weighted}: $A^{(\ell)} = \sum_{h=1}^H w_h A^{(\ell)}_h$ with learned or value-dependent weights
\end{enumerate}

\subsection{Attention Flow via Max-Flow}

An alternative to rollout is computing attention flow using network flow algorithms:

\begin{definition}[Attention Flow]
Model the transformer as a flow network where:
\begin{itemize}
    \item Nodes represent (layer, position) pairs
    \item Edge capacities are attention weights
    \item Add source/sink nodes for input/output positions
\end{itemize}
The attention flow from input position $j$ to output position $i$ is the maximum flow from the source connected to $(0, j)$ to the sink connected to $(L, i)$.
\end{definition}

Attention flow satisfies conservation properties that rollout does not, providing a more principled measure of information propagation.

% =============================================================================
\section{Uncertainty Quantification via Semantic Entropy}
% =============================================================================

\subsection{The Problem with Token-Level Uncertainty}

Standard uncertainty measures for language models operate at the token level:
\begin{equation}
H(x_t \mid x_{<t}) = -\sum_{v \in \V} p(x_t = v \mid x_{<t}) \log p(x_t = v \mid x_{<t})
\end{equation}

However, this conflates linguistic uncertainty (many ways to phrase the same meaning) with semantic uncertainty (uncertainty about the underlying facts).

\subsection{Semantic Equivalence Classes}

The key insight is to partition outputs into semantic equivalence classes.

\begin{definition}[Semantic Equivalence]
Two sequences $s_1, s_2$ are semantically equivalent, denoted $s_1 \sim s_2$, if they express the same meaning. This induces a partition of the output space into equivalence classes $[s] = \{s' : s' \sim s\}$.
\end{definition}

In practice, semantic equivalence is approximated using:
\begin{enumerate}
    \item \textbf{Bidirectional entailment}: $s_1 \sim s_2$ if an NLI model predicts $s_1 \Rightarrow s_2$ and $s_2 \Rightarrow s_1$
    \item \textbf{Embedding similarity}: $s_1 \sim s_2$ if $\text{sim}(e(s_1), e(s_2)) > \tau$ for sentence embeddings $e$
\end{enumerate}

\subsection{Semantic Entropy}

\begin{definition}[Semantic Entropy]
Given a model $p$ and context $c$, the semantic entropy is:
\begin{equation}
\text{SE}(c) = -\sum_{[s] \in \mathcal{C}/{\sim}} p([s] \mid c) \log p([s] \mid c)
\end{equation}
where $\mathcal{C}/{\sim}$ denotes the set of semantic equivalence classes and:
\begin{equation}
p([s] \mid c) = \sum_{s' \in [s]} p(s' \mid c)
\end{equation}
\end{definition}

\subsubsection{Estimation Procedure}

Since exact computation is intractable, semantic entropy is estimated via sampling:

\begin{enumerate}
    \item Generate $K$ samples $s_1, \ldots, s_K$ from $p(\cdot \mid c)$
    \item Cluster samples into semantic equivalence classes using pairwise entailment checks
    \item Estimate class probabilities: $\hat{p}([s]) = \frac{1}{K}\sum_{i=1}^K \mathbf{1}[s_i \in [s]]$
    \item Compute entropy: $\widehat{\text{SE}}(c) = -\sum_{[s]} \hat{p}([s]) \log \hat{p}([s])$
\end{enumerate}

\subsection{Properties of Semantic Entropy}

\begin{proposition}[Bounds]
For any semantic clustering:
\begin{equation}
\text{SE}(c) \leq H(S \mid c)
\end{equation}
where $H(S \mid c)$ is the standard sequence entropy. Equality holds when each equivalence class contains exactly one sequence.
\end{proposition}

\begin{proposition}[Factuality Correlation]
In empirical studies, semantic entropy correlates more strongly with factual accuracy than:
\begin{itemize}
    \item Token-level entropy
    \item Sequence probability
    \item Sample consistency measures
\end{itemize}
This is because semantic entropy captures uncertainty about meaning rather than phrasing.
\end{proposition}

% =============================================================================
\section{Hallucination Detection Methods}
% =============================================================================

\subsection{SelfCheckGPT: Black-Box Detection via Sampling Consistency}

SelfCheckGPT detects hallucinations without access to model internals by exploiting the observation that factual information tends to be reproduced consistently across samples, while hallucinated content varies.

\subsubsection{Core Principle}

Given a response $r$ generated from prompt $p$, generate additional samples $s_1, \ldots, s_K$ from the same prompt. If $r$ contains factual information, the samples should be consistent; if $r$ contains hallucinations, the samples will likely contradict each other.

\subsubsection{Consistency Metrics}

Several metrics quantify consistency between the original response and samples:

\paragraph{BERTScore Variant}
For each sentence $r_i$ in the response:
\begin{equation}
\text{Score}_{\text{BERT}}(r_i) = 1 - \frac{1}{K}\sum_{k=1}^K \text{BERTScore}(r_i, s_k)
\end{equation}
Higher scores indicate potential hallucination.

\paragraph{NLI-Based}
Using a natural language inference model:
\begin{equation}
\text{Score}_{\text{NLI}}(r_i) = \frac{1}{K}\sum_{k=1}^K \mathbf{1}[\text{NLI}(s_k, r_i) = \text{contradiction}]
\end{equation}

\paragraph{Question-Answering}
Generate questions about $r_i$, answer them using each sample $s_k$, and compare:
\begin{equation}
\text{Score}_{\text{QA}}(r_i) = 1 - \frac{1}{|\mathcal{Q}|}\sum_{q \in \mathcal{Q}} \frac{1}{K}\sum_{k=1}^K \text{F1}(\text{Answer}(q, r_i), \text{Answer}(q, s_k))
\end{equation}

\paragraph{N-gram Consistency}
\begin{equation}
\text{Score}_{n\text{-gram}}(r_i) = 1 - \frac{1}{K}\sum_{k=1}^K \frac{|\text{ngrams}(r_i) \cap \text{ngrams}(s_k)|}{|\text{ngrams}(r_i)|}
\end{equation}

\subsubsection{Aggregation}

Sentence-level scores are aggregated to passage-level:
\begin{equation}
\text{Score}(r) = \text{Aggregate}(\text{Score}(r_1), \ldots, \text{Score}(r_m))
\end{equation}
where $\text{Aggregate}$ can be mean, max, or a weighted combination.

\subsection{DoLa: Decoding by Contrasting Layers}

DoLa improves factuality by exploiting the observation that factual knowledge emerges in later transformer layers.

\subsubsection{Motivation: Knowledge Localization}

Empirical studies show that transformer layers encode different types of information:
\begin{itemize}
    \item \textbf{Early layers}: Syntactic patterns, surface-level features
    \item \textbf{Middle layers}: Semantic relationships
    \item \textbf{Late layers}: Factual knowledge, world knowledge
\end{itemize}

When predicting factual content, the probability distribution changes significantly in the final layers as factual knowledge is injected.

\subsubsection{Layer Contrasting}

Let $q_j(x_t \mid x_{<t})$ denote the next-token distribution obtained by applying the language model head to the hidden state at layer $j$:
\begin{equation}
q_j(x_t \mid x_{<t}) = \softmax(\phi(h_t^{(j)}))
\end{equation}

DoLa computes the final distribution by contrasting a ``mature'' layer $N$ (final layer) with a ``premature'' layer $M$:
\begin{equation}
\hat{p}(x_t \mid x_{<t}) = \softmax\left(\log q_N(x_t) - \log q_M(x_t)\right)
\end{equation}

\subsubsection{Dynamic Layer Selection}

The premature layer $M$ is selected dynamically at each decoding step to maximize the distributional divergence:
\begin{equation}
M = \argmax_{j \in \mathcal{J}} \JSD(q_N(\cdot \mid x_{<t}) \| q_j(\cdot \mid x_{<t}))
\end{equation}
where $\mathcal{J}$ is a candidate set of early layers and JSD is the Jensen-Shannon divergence.

The intuition: when predicting tokens requiring factual knowledge, the distribution changes significantly between early and late layers. Selecting the layer with maximum divergence captures this ``knowledge injection'' point.

\subsubsection{Adaptive Plausibility Constraint}

To prevent implausible tokens from receiving high scores after contrasting:
\begin{equation}
\hat{p}(x_t) \propto \begin{cases}
\frac{q_N(x_t)}{q_M(x_t)} & \text{if } q_N(x_t) \geq \alpha \max_w q_N(w) \\
0 & \text{otherwise}
\end{cases}
\end{equation}

This ensures only tokens that are plausible according to the mature layer can be selected.

\subsubsection{Experimental Results}

On TruthfulQA, DoLa improves LLaMA models by 12-17\% absolute on truthfulness metrics. Key findings:

\begin{center}
\begin{tabular}{lcc}
\hline
Task Type & Optimal Premature Layers & Interpretation \\
\hline
Short factual (TruthfulQA) & Higher layers (e.g., [16,32) for 7B) & Facts in final layers \\
Long paragraph (FACTOR) & Lower layers (e.g., [0,16) for 7B) & Less layer-specific \\
Chain-of-thought (GSM8K) & Lower layers & Reasoning distributed \\
\hline
\end{tabular}
\end{center}

The efficiency overhead is minimal: 1-8\% increase in latency due to early exit computations.

% =============================================================================
\section{Unified Theoretical Framework}
% =============================================================================

\subsection{Connecting Information Theory to Detection}

The methods described above can be unified through the lens of bounded information theory:

\begin{enumerate}
    \item \textbf{V-Information}: The model's ``usable information'' about facts depends on its architecture ($\F$). Hallucinations occur when the model lacks sufficient $\F$-information about the query.

    \item \textbf{Attention Flow}: Information from context propagates through attention. Low attention flow from relevant context to generated tokens suggests the model is not using contextual information (potential hallucination).

    \item \textbf{Semantic Entropy}: High semantic entropy indicates the model's $\F$-entropy about the true answer is high---multiple semantically distinct answers are plausible given the model's bounded capacity.

    \item \textbf{SelfCheckGPT}: Sampling consistency measures empirical variance in the model's outputs. High variance suggests the model lacks stable $\F$-information.

    \item \textbf{DoLa}: Layer contrasting identifies where factual $\F$-information is encoded. By emphasizing late-layer knowledge, we increase reliance on factual information.
\end{enumerate}

\subsection{The Authority-Trust Framework}

These insights motivate a unified detection framework:

\begin{equation}
\text{Uncertainty}(t) = 1 - \text{Authority}(t)
\end{equation}

where Authority combines:
\begin{itemize}
    \item \textbf{Information Flow}: How much context information reaches token $t$
    \item \textbf{Layer Agreement}: Whether early and late layers agree (DoLa-inspired)
    \item \textbf{Semantic Stability}: Whether the prediction is semantically consistent (SelfCheck-inspired)
\end{itemize}

Formally:
\begin{equation}
\text{Authority}(t) = \text{Gate}(t) \cdot \text{Flow}(t) + (1 - \text{Gate}(t)) \cdot \text{Trust}(t)
\end{equation}
\begin{equation}
\text{Trust}(t) = 1 - \text{Dispersion}(t) \cdot (1 + \lambda \cdot \text{Varentropy}(t))
\end{equation}

where:
\begin{itemize}
    \item $\text{Flow}(t)$: Attention-based information flow from context
    \item $\text{Gate}(t)$: MLP-attention agreement (layer-contrasting inspired)
    \item $\text{Dispersion}(t)$: Semantic dispersion across top-k predictions
    \item $\text{Varentropy}(t)$: Variance of log-probabilities (confidence stability)
\end{itemize}

% =============================================================================
\section{Experimental Benchmarks and Metrics}
% =============================================================================

\subsection{Standard Benchmarks}

\paragraph{TruthfulQA} Tests models' tendency to produce truthful answers versus common misconceptions. Metrics include MC1 (single correct), MC2/MC3 (multiple correct), and human-evaluated truthfulness/informativeness scores.

\paragraph{HaluEval} Specifically designed for hallucination detection with QA, summarization, and dialogue tasks. Each example includes human-annotated hallucination labels.

\paragraph{FACTOR} Long-form factuality benchmark with News and Wiki domains. Tests ability to identify correct vs incorrect paragraph completions.

\paragraph{RAGTruth} Retrieval-augmented generation benchmark testing faithfulness to retrieved context.

\subsection{Evaluation Metrics}

\begin{itemize}
    \item \textbf{AUROC}: Area under ROC curve for hallucination detection as binary classification
    \item \textbf{AUPRC}: Area under precision-recall curve (important for class imbalance)
    \item \textbf{ECE}: Expected calibration error between confidence and accuracy
    \item \textbf{Spearman/Pearson}: Correlation between uncertainty scores and error rates
\end{itemize}

% =============================================================================
\section{Discussion and Future Directions}
% =============================================================================

\subsection{Key Insights}

\begin{enumerate}
    \item \textbf{Bounded computation matters}: Traditional information theory fails to capture what models can actually learn and use. V-information and epiplexity provide better frameworks.

    \item \textbf{Layer structure encodes knowledge type}: Different types of information (syntax, semantics, facts) localize to different layers. This structure can be exploited for detection.

    \item \textbf{Semantic vs linguistic uncertainty}: Distinguishing uncertainty about meaning from uncertainty about phrasing is crucial for reliable uncertainty quantification.

    \item \textbf{Consistency as a signal}: Both sampling-based (SelfCheckGPT) and layer-based (DoLa) methods exploit consistency as a hallucination signal.
\end{enumerate}

\subsection{Open Problems}

\begin{enumerate}
    \item \textbf{Calibration}: How to calibrate uncertainty estimates across different query types and domains?

    \item \textbf{Efficiency}: How to achieve robust detection without multiple samples or expensive layer computations?

    \item \textbf{Attribution}: Beyond detection, how to identify the source of hallucinations (parametric knowledge vs context misuse)?

    \item \textbf{Mitigation}: How to use detection signals to prevent hallucinations during generation?
\end{enumerate}

% =============================================================================
\section{Conclusion}
% =============================================================================

This document has presented a unified view of the theoretical foundations underlying hallucination detection in large language models. By connecting information theory under computational constraints, attention flow analysis, semantic uncertainty quantification, and practical detection methods, we provide a comprehensive framework for understanding and addressing the hallucination problem.

The key unifying insight is that hallucinations arise from limitations in the model's bounded information processing capacity. Detection methods that account for these limitations---by measuring information flow, layer-wise knowledge localization, and semantic consistency---outperform naive approaches based on standard uncertainty measures.

% =============================================================================
% References (Synthesized)
% =============================================================================

\section*{Source Papers}

\begin{enumerate}
    \item Xu, Y., Cao, H., Du, Y., Shao, H., Gao, Y., \& Yang, Y. (2020). A Theory of Usable Information under Computational Constraints. \textit{ICLR 2020}.

    \item Abnar, S., \& Zuidema, W. (2020). Quantifying Attention Flow in Transformers. \textit{ACL 2020}.

    \item Kuhn, L., Gal, Y., \& Farquhar, S. (2023). Semantic Uncertainty: Linguistic Invariances for Uncertainty Estimation in Natural Language Generation. \textit{ICLR 2023}.

    \item Manakul, P., Liusie, A., \& Gales, M. J. F. (2023). SelfCheckGPT: Zero-Resource Black-Box Hallucination Detection for Generative Large Language Models. \textit{EMNLP 2023}.

    \item Chuang, Y.-S., Xie, Y., Luo, H., Kim, Y., Glass, J., \& He, P. (2024). DoLa: Decoding by Contrasting Layers Improves Factuality in Large Language Models. \textit{ICLR 2024}.

    \item Finzi, M., Qiu, S., Jiang, Y., Izmailov, P., Kolter, J. Z., \& Wilson, A. G. (2025). From Entropy to Epiplexity: Rethinking Information for Computationally Bounded Intelligence. \textit{JMLR}.
\end{enumerate}

\end{document}
